% ID: SEL-D0041-OB-E03
% ESTADO: cotejada. Segunda pasada visual y borrador compilado revisados.
% Fuente: Selectividad/matematicas_aplicadas/Ua6MATEMÁTICAS APLICADAS examen.pdf
% SHA-256: b7541cc45d3e96297c29037fe17a562b7daa9cd0ee2a95c2b5e9454ad9b90f14
% Página: 2; opción B; ejercicio original 3.

Una enfermedad puede estar provocada por solo una de estas tres causas: A, B o C. La probabilidad de que la causa sea A es 0.3, la de que sea B es 0.2 y la de que sea C es 0.5. El tratamiento de esta enfermedad requiere hospitalización en el 20\% de los casos si está provocada por A, en el 55\% si la causa es B y en el 10\% si la causa es C.
\begin{enumerate}[label=\alph*)]
\item \textbf{(1.5 puntos)} \textquestiondown{}Cuál es la probabilidad de que un enfermo con la citada enfermedad no necesite hospitalización?
\item \textbf{(1 punto)} Si un enfermo está hospitalizado debido a esta enfermedad, \textquestiondown{}cuál es la probabilidad de que la causa haya sido A?
\end{enumerate}
